\section{Introduction}
\label{sec:intro}

Reconstructing high-fidelity 3D scenes for photorealistic novel view synthesis remains a core challenge in computer vision. While Neural Radiance Fields (NeRF)~\cite{mildenhall2021nerf} and its variants~\cite{barron2021mipnerf,Barron_2022_CVPR,barron2023zipnerf} achieve remarkable visual quality, their implicit representation incurs slow rendering. Recently, 3D Gaussian Splatting (3DGS)~\cite{kerbl20233d} introduces an explicit and efficient alternative, delivering NeRF-level fidelity with real-time performance, and a series of its extensions~\cite{Yu2023MipSplatting,sabourgoli2024spotlesssplats,seo2025flod} rapidly establishing a new state-of-the-art for high-quality, efficient scene reconstruction.
%-------------------------------------------------------------------------

With the success of these methods in dense-view scenarios, the research frontier is now shifting toward the more challenging sparse-input setting. In this regime, the scarcity of viewpoints often leads to irregular Gaussian distributions, resulting in geometric ambiguities and severe rendering artifacts. 
FSGS~\cite{zhu2024fsgs} accelerates early-stage densification to compensate for the insufficient initial Gaussians, while DropGaussian~\cite{Park_2025_CVPR} adopts a dropout-based strategy to regularize the Gaussian distributions in occluded regions. 
However, the absence of comprehensive guidance still leaves the Gaussian field imperfectly optimized, resulting in inaccurate and uneven distributions — with insufficient Gaussians in background regions leading to blurriness, and inadequate supervision over high-frequency details causing misaligned or misplaced Gaussians.
This motivates us to introduce a multi-level framework that globally-to-locally coherent guidance to more comprehensively direct the formation of accurate Gaussian distributions.

Sparse-view settings often cause Gaussians to fall outside the field of view, receiving limited gradient feedback and overfitting to a few training images~\cite{Park_2025_CVPR}.
As shown in Fig.~\ref{fig:new_motivation}, adding more views improves gradient coverage and alleviates this issue.
Inspired by this, we synthesize pseudo-labels that, together with real views, provide dense image-level supervision to regularize Gaussian distributions.
This enhanced supervision enriches gradient propagation across the scene, leading to more uniform Gaussian distributions and improved reconstruction fidelity under sparse inputs, while providing richer spatial structural cues for fine-grained guidance at the feature level.

To further refine and regularize the Gaussian distributions in high-frequency regions, we introduce the \textbf{F}eature-\textbf{A}daptive \textbf{D}ensification and \textbf{P}runing (\textbf{FADP}) at the feature level.
FADP increases Gaussian density along edge-aware regions to capture high-frequency details, prunes redundant splats within homogeneous patches to avoid over-saturation, and adds new Gaussians in sparse background areas to improve coverage.
The refined Gaussian distributions from FADP provide a stable foundation for the parameter-level, enabling more reliable co-pruning based on geometric consistency.
\begin{figure}[t]
  \centering
  \includegraphics[width=0.95\columnwidth]{imgs/Fig1_Motivation.pdf}
%   \vspace{-7pt}

  \caption{ 
  \textbf{Motivation of HeroGS.} 
  (Left) Image-level pseudo-labels bridge the gap between sparse and dense supervision, yielding a more complete Gaussian distributions. 
  (Right) FADP and CPG refine inaccurate Gaussians by enhancing distributions closer to ground-truth geometry and pruning inconsistencies. 
 % in SideFormer
 }
  \vspace{-10pt}
  \label{fig:new_motivation}
\end{figure}

%In Fig.~\ref{fig:vis_detail_depth}C, it clearly illustrates the significant deterioration in reconstruction quality for 3DGS and its derivative methods as input views decrease. In contrast, our proposed framework exhibits enhanced robustness under sparse-view conditions, effectively mitigating the performance drop and maintaining higher fidelity in both geometry and appearance (Fig.~\ref{fig:vis_detail_depth}A). Quantitatively, the proposed method outperforms state-of-the-art approaches across multiple evaluation metrics—particularly under extremely sparse inputs—underscoring its superior capability in high-quality reconstruction from minimal data. 
At the parameter level, we introduce the \textbf{C}o-\textbf{P}runed \textbf{G}eometry Consistency (\textbf{CPG}), which eliminates abnormal or inconsistent Gaussian distributions via a co-pruning mechanism combined with a post-freeze behavior strategy.
Co-Pruning leverages the self-consistency of the Gaussian field to jointly evaluate and filter Gaussian parameters, retaining only those that exhibit stable and geometrically consistent distributions.
This process effectively preserves robust Gaussians while eliminating unstable or redundant ones, leading to a cleaner and more reliable scene representation, as illustrated in the right part of Fig.~\ref{fig:new_motivation}.
The whole pipeline constitutes a hierarchical guidance strategy that collaboratively optimizes the Gaussian distributions across levels, enhancing its global consistency and fidelity.
Our contributions can be summarized as follows:

\begin{itemize}
%%%%%你这一个点写太多了和第二个点对比太鲜明了，一个超长一个超短，你分三个四个点写都行啊
   \item  We propose \textbf{HeroGS}, a hierarchical guidance framework that enables compact and high-fidelity 3D reconstructions under sparse-view conditions.
   \item At image level, pseudo-labels are generated to promote more accurate Gaussian distributions. At  feature and parameter levels, the proposed \textbf{FADP} and \textbf{CPG} further refine Gaussian field by enhancing feature-aware density and enforcing geometric consistency, respectively.


   \item Extensive qualitative and quantitative experiments on real-world datasets, including large-scale scenes and various training views, demonstrate that our method significantly outperforms state-of-the-art baselines.
\end{itemize}


\section{Related Work}
\begin{figure*}[t]
  \centering
  \includegraphics[width=0.95\textwidth]{imgs/Fig2_HeroGS.pdf}
%   \vspace{-7pt}
  \caption{ \textbf{HeroGS Overview.} Initialized from SfM, our framework operates across three levels.
(1) \textbf{Image-level Guidance:} A set of intermediate RGB frames is synthesized from sparse inputs, offering pseudo-dense guidance that globally regularizes the Gaussian distributions and, at the feature level, manifests as patch-based Gaussian numbers $\mathcal{C}$.
%and jointly optimize their poses together with 3D Gaussian field.
(2) \textbf{Feature-level Refinement:} The Feature-Adaptive Densification and Pruning (FADP) leverages edge- and patch-aware features from training views to enhance high-frequency and background regions, while suppressing redundant Gaussians and dilivering finer details for next level.
(3) \textbf{Parameter-level Consistency:} The Co-Pruned Geometry Consistency (CPG) employs auxiliary Gaussian fields with partially frozen parameters to perform co-pruning, eliminating geometrically inconsistent splats.
These levels form a hierarchical guidance with interconnections (dashed lines) that jointly constrains and optimizes the Gaussian field for improved structural fidelity.
  %This paradigm introduces a comprehensive Pseudo-Label framework composed of multiple interacting modules. We first manufacture basic but noisy pseudo-labels by interpolating intermediate RGB frames to provide auxiliary supervision. Here, we jointly optimize the Gaussian Splatting representation(initialized by SFM) and the poses of the interpolated views using a combination of rendered loss $\mathcal{L}_{\mathrm{g}}$ and RGB loss $\mathcal{L}_{\mathrm{r}}$. Second, a Feature-Adaptive Densification and Pruning (FADP) module rectifies the pseudo-label-based training process by selectively densifying and pruning Gaussians based on features extracted from training views. To purge erroneous pseudo-label influence, two auxiliary Gaussian fields with partially frozen parameters are introduced for co-pruning to suppress inconsistent Gaussians. Together, these components form a unified pseudo-label supervision paradigm that enhances geometric consistency and improves rendering quality under sparse input conditions.
 % in SideFormer
 }
  \vspace{-5pt}
  \label{fig:framework}
\end{figure*}
%%%%%不要用\para这个符号，两段之间空的太多了，用\textbf加粗一下就行
\textbf{Novel View Synthesis.}
\textbf{Neural Radiance Fields (NeRF)}~\cite{mildenhall2021nerf, barron2022mip, bao2023and} and \textbf{3D Gaussian Splatting (3DGS)}~\cite{kerbl20233d} have emerged as two prominent paradigms for high-fidelity novel view synthesis. NeRF-based methods have achieved impressive results in producing photorealistic novel views, especially when abundant views (often hundreds) are available. Meanwhile, 3DGS represents scenes explicitly as a set of 3D Gaussian primitives, which enables faster convergence and real-time rendering performance through GPU-friendly splatting operations. Recent works such as MiniSplatting~\cite{kaiser2024minisplatting} focus on optimizing 3DGS by introducing memory-efficient representation and lightweight pipeline, while Stop-the-Pop~\cite{radl2024stopthepop} improves rendering realism and robustness by a novel hierarchical rasterization approach. Beyond these, Deformable 3DGaussians~\cite{yang2023deformable3dgs} extends 3DGS to efficiently reconstruct and render dynamic scenes by learning a canonical static scene and a time-dependent deformation field for Gaussians. However, both categories are fundamentally data-hungry and show performance degradation under sparse input settings, where their reliance on dense photometric supervision becomes a limiting factor.

\textbf{Novel View Synthesis with sparse views.}
%%High-fidelity 3D reconstruction and novel view synthesis techniques, such as Neural Radiance Fields (NeRF) ~\cite{wang2023sparsenerf, mildenhall2021nerf, barron2022mip, bao2023and} and 3D Gaussian Splatting (3DGS) ~\cite{ kerbl20233d}, have made significant progress in generating realistic images. However, these methods typically require a large number of input images (hundreds) to learn a scene, which limits their practical application.
To address the challenges of sparse views, existing methods based on NeRF are mainly divided into two categories: (1) regularization-based methods, which apply techniques such as depth smoothness~\cite{ niemeyer2022regnerf}, occlusion regularization~\cite{yang2023freenerf}, or frequency control to existing sparse data to prevent overfitting; (2) methods incorporating external priors like using pre-trained models to generate depth maps~\cite{chen2021mvsnerf, fan2024instantsplat, jang2021codenerf, yu2021pixelnerf} or feature extractors~\cite{jain2021putting, kwak2023geconerf} to enhance geometric and visual consistency. For 3DGS specifically, some methods have been proposed, such as DNGaussian~\cite{li2024dngaussian} introducing depth regularization, SparseGS~\cite{xiong2024sparsegs} combining depth and diffusion regularization, FSGS~\cite{zhu2024fsgs} increasing point cloud density through Gaussian unpooling, and CoR-GS~\cite{zhang2024cor} and CoherentGS~\cite{paliwal2024coherentgs} utilizing multi-view consistency or monocular depth for optimization. However, existing sparse-view 3D reconstruction methods generally still suffer from severe overfitting, indicating a need for further research to improve their robustness in practical applications.

%-------------------------------------------------------------------------














